% RecSys Odyssey repository-backed lecture note.
% Add figures with: \coursefigure{figures/name.svg}{Caption}
\section{Estimate counterfactual value}

\begin{abstract}
Propensity weighting asks how outcomes would look under a different exposure policy.
\end{abstract}

\subsection{Concept}
Inverse Propensity Scoring weights an observed reward by the inverse probability that the logging policy chose that action. Rarely shown actions then carry more information. Self-normalized and clipped variants control variance; doubly robust estimators add a reward model. None can repair missing overlap: if the old policy never exposed an action, the log contains no evidence about its outcome.

\begin{equation*}
r̂_IPS = [π_target(a|x) / π_log(a|x)] \cdot r
\end{equation*}

\subsection{Key terms}
\begin{itemize}
\item \textbf{Propensity.} The logging policy’s probability of selecting an observed action.
\item \textbf{Counterfactual.} The unobserved outcome under another possible action or policy.
\item \textbf{Weight clipping.} Limiting extreme propensity weights to reduce estimator variance.
\end{itemize}
